%%%%%%%%%%%%%%%%%%%%%%%%%%%%%%%%%%%%%%%%%%%%%%%%%%%%%%%%%%%%%%%%%
%%% %
%%% % weiiszablon.tex
%%% % The Faculty of Electrical and Computer Engineering
%%% % Rzeszow University Of Technology diploma thesis Template
%%% % Szablon pracy dyplomowej Wydziału Elektrotechniki 
%%% % i Informatyki PRz
%%% % January, 2024
%%%%%%%%%%%%%%%%%%%%%%%%%%%%%%%%%%%%%%%%%%%%%%%%%%%%%%%%%%%%%%%%%

\documentclass[12pt,twoside,openright]{article}

\usepackage{weiiszablon}

\author{Kacper Hyliński}

% np. EF-123456, EN-654321, ...
\studentID{177082}

\title{Implementacja algorytmu DQN do budowania strategii agenta gry komputerowej}
\titleEN{Implementation of the DQN algorithm for building a computer game agent's strategy}


%%% wybierz rodzaj pracy wpisując jeden z poniższych numerów: ...
% 1 = inżynierska	% BSc
% 2 = magisterska	% MSc
% 3 = doktorska		% PhD
% 4 = praca inżynierska
%%% na miejsce zera w linijce poniżej
\newcommand{\rodzajPracyNo}{1}


%%% promotor
\supervisor{(dr hab. inż.) Roman Zajdel (prof. PRz)}
%% przykład: dr hab. inż. Józef Nowak, prof. PRz

%%% promotor ze stopniami naukowymi po angielsku
\supervisorEN{Roman Zajdel, DSc, Eng., Associate Professor}

\abstract{Treść streszczenia po polsku}
\abstractEN{Treść streszczenia po angielsku}

\keywords{max. 5 słów kluczowych w 2 wierszach, oddzielanych przecinkami}
\keywordsEN{max. five keywords in English}


\begin{document}

% strona tytułowa
\maketitle

\cleardoublepage

% spis treści
\tableofcontents

\cleardoublepage

% !TEX root = ../praca_inzynierska.tex
\section*{Wykaz symboli, oznaczeń i skrótów}
\addcontentsline{toc}{section}{Wykaz symboli, oznaczeń i skrótów}
% !TEX root = ../praca_inzynierska.tex
\section{Wstęp}
\label{sec:wstep}

\subsection{Motywacja i zarys problemu}
\label{subsec:motywacja}

Gry komputerowe od lat stanowią jedno z najważniejszych środowisk badawczych sztucznej inteligencji. Mają jasno określone reguły, mierzalny wynik i pozwalają przeprowadzić tysiące rozgrywek w czasie, w którym człowiek rozegrałby zaledwie kilka. Te cechy sprawiają, że gry dobrze nadają się do badania metod uczenia ze wzmocnieniem, w których agent uczy się strategii działania wyłącznie na podstawie nagród otrzymywanych za swoje decyzje \cite{sutton2018}.

Przełomem w tej dziedzinie było opublikowanie w 2015 roku algorytmu Deep Q-Network (DQN). Agent DQN, otrzymując na wejściu jedynie obraz ekranu i wynik gry, nauczył się grać w 49 gier z konsoli Atari 2600 i w wielu z nich osiągnął poziom porównywalny z profesjonalnym graczem \cite{mnih2015}. Kluczowym pomysłem było połączenie klasycznego algorytmu Q-learning \cite{watkins1992} z głęboką siecią neuronową, która przybliża wartości akcji w stanach zbyt licznych, aby można je było zapisać w tablicy. Kolejne prace pokazały, że uczenie ze wzmocnieniem pozwala tworzyć agentów pokonujących mistrzów świata w wywodzącej się z Chin grze Go \cite{silver2016} oraz osiągających poziom najlepszych graczy w złożonej grze strategicznej czasu rzeczywistego StarCraft II \cite{vinyals2019}.

Osiągnięcia te wymagały jednak ogromnych zasobów obliczeniowych, niedostępnych poza dużymi laboratoriami badawczymi. Aby samodzielnie zbadać działanie algorytmu DQN, potrzebne jest środowisko, które z jednej strony stawia agentowi rzeczywiste wyzwania, a z drugiej pozwala przeprowadzić trening na komputerze osobistym. Takim środowiskiem jest gra Snake. Jej reguły są proste: wąż porusza się po planszy, zjada pożywienie i rośnie, a gra kończy się w chwili zderzenia ze ścianą lub z własnym ciałem. Mimo to gra wymaga od agenta planowania kilku ruchów naprzód, unikania zagrożeń, które zmieniają się wraz z długością węża, oraz łączenia celu krótkoterminowego (zdobycie pożywienia) z długoterminowym (przetrwanie).

Dodatkową motywacją do podjęcia tematu jest możliwość praktycznego sprawdzenia rozwiązań zaproponowanych w literaturze jako ulepszenia podstawowego algorytmu DQN, takich jak Double DQN \cite{vanhasselt2016}, oraz porównania własnej implementacji z dostępnymi rozwiązaniami bibliotecznymi.

Grę Snake można rozwiązywać również bez uczenia, metodami klasycznego przeszukiwania. Algorytm A* \cite{hart1968} wyznacza najkrótszą ścieżkę od głowy węża do pożywienia, korzystając z pełnej wiedzy o planszy i heurystyki odległości. Zestawienie go z agentem DQN, który nie zna reguł gry i uczy się wyłącznie na podstawie nagród, pozwala ocenić, co zyskuje się, a co traci, zastępując zaprogramowaną strategię strategią wyuczoną.
\subsection{Struktura pracy}
\label{subsec:struktura}

Praca składa się z sześciu rozdziałów. W rozdziale \ref{sec:rl-gry} przedstawiono podstawy teoretyczne uczenia ze wzmocnieniem oraz omówiono wykorzystanie tych metod w grach komputerowych. Rozdział \ref{sec:przeglad} zawiera przegląd algorytmów uczenia ze wzmocnieniem, klasycznych metod przeszukiwania i bibliotek programistycznych, zakończony uzasadnieniem wyboru algorytmu DQN. W rozdziale \ref{sec:projekt} opisano projekt systemu: środowisko gry, reprezentację stanu, funkcję nagrody i architekturę aplikacji. Rozdział \ref{sec:implementacja} poświęcono implementacji środowiska, sieci neuronowej, agenta i procesu uczenia.

% Wyniki przeprowadzonych eksperymentów przedstawiono w rozdziale \ref{sec:wyniki}.
Ostatni rozdział zawiera podsumowanie pracy i wnioski.

\subsection{Cel i zakres pracy}
\label{subsec:cel}

Celem pracy było przedstawienie wykorzystania algorytmów uczenia ze wzmocnieniem w grach komputerowych, dokonanie przeglądu tych algorytmów oraz dostępnych rozwiązań bibliotecznych, a także implementacja algorytmu DQN do budowania strategii agenta gry Snake w języku Python, oraz porównanie go z agentem wykorzystującym algorytm przeszukiwania A*. Zakres pracy obejmował opracowanie środowiska gry w wersji graficznej i uproszczonej, przyspieszonej na potrzeby treningu, zaprojektowanie reprezentacji stanu i funkcji nagrody, implementację agenta DQN w wariancie Double DQN z wykorzystaniem biblioteki PyTorch oraz przygotowanie procesu uczenia z możliwością wykorzystania procesora graficznego.
% Dopisać po decyzji promotora: „oraz ocenę skuteczności wytrenowanego agenta”.
\input{rozdzialy/02-teoria}
\input{rozdzialy/03-projekt}
\input{rozdzialy/04-implementacja}
% \input{rozdzialy/05-wyniki}   % czeka na decyzję promotora
\input{rozdzialy/06-podsumowanie}

\section*{Wykaz symboli, oznaczeń i skrótów (opcjonalny)}
\addcontentsline{toc}{section}{Wykaz symboli, oznaczeń i skrótów (opcjonalny)}%

1 $\div$ 2 stron wykaz ważniejszych symboli i oznaczeń (jeśli jest potrzebny).


\clearpage

\addcontentsline{toc}{section}{Literatura}
\bibliographystyle{unsrt}
\bibliography{literatura}

\clearpage

\makesummary

\end{document} 
